\section{A tensor product of \texorpdfstring{\cstars{}s}{σC*-algebras}}\label{sec:tensor}

Our next aim is to ensure that $\SigmaC$ and $\SigmaCcpu$ have a tensor product.
From now on,
\[
	\otimes\colon \Calg \times \Calg \to \Calg
\]
denotes the \newterm{maximal tensor product}, satisfying the following universal property: for every $\phi \colon \A \to \C$ and $\psi \colon \B \to \C$ with commuting images, there exists a unique $(\phi,\psi)\colon \A \otimes \B \to \C$ such that $(\phi,\psi)(a\otimes b)= \phi(a)\psi(b)$.
The resulting triple $(\Calg,\otimes,\bbC)$ is a symmetric monoidal category.
Throughout, we shall work exclusively with this monoidal structure.\footnote{In particular, the minimal tensor product will not be discussed.}

This tensor is known to extend to cpu maps with the same universal property, so in particular also $(\Calgcpu,\otimes ,\bbC)$ is a symmetric monoidal category.
For more details, the reader may refer to \cite{brownozawa}, particularly Section 3.3 and Exercise 3.5.1.


To equip $\SigmaC$ with a tensor product, we rely on its monadicity over $\Calg$ (\cref{thm:U_monadic}). 
The construction follows the well-known description of tensor products of algebras over monoidal monads, first investigated by Kock in the early 1970s~\cite{Kock1970}.
We employ results from \cite{seal2013tensors}, where this machinery was generalized to monoidal categories that are not necessarily closed, as is the case for $(\Calg,\otimes,\bbC)$.

Since monoidality of the monad is necessary, we focus on this first.
To prove it, we only rely on well-established results, so the reader is not required to know the precise definition.
\begin{proposition}\label{prop:monoidalmonad}
	The monad $T=UF$ obtained from the adjunction of \autoref{thm:adj_sigma} is a symmetric monoidal monad in the sense of \cite{Kock1970}.
\end{proposition}

\begin{proof}
	The Kleisli category of $T$ is given by Pedersen--Baire envelopes with $\sigma$-homomorphisms by direct check.
	This was proven to be symmetric monoidal in \cite[Proposition 5.3.2]{fritz2025categories}, and so $T$ is a symmetric monoidal monad by \cite[Proposition 1.2.2]{seal2013tensors}.
\end{proof}
Monoidality of $T$ is witnessed by the maps
\[ \theta_{\A,\B}\colon T\A \otimes T\B \to T(\A\otimes \B),\]
uniquely determined by the universal property of the maximal tensor product applied to the $\sigma$-homomorphisms $T(\A)\to T(\A\otimes B)$ and $T(\B)\to T(\A\otimes \B)$, which are uniquely obtained by
\begin{equation*}%\label{eq:description_theta}
	\A\owns a\mapsto \eta_{A\otimes B}(a\otimes 1) \in T(\A\otimes B)\qquad \text{and}\qquad \B \owns b\mapsto \eta_{A\otimes B}(1\otimes b) \in T(\A\otimes B).
\end{equation*}

% It is worth noting that, for order convergent sequences $(\eta_{\A} (a_n))_n\subseteq T\A$ and $(\eta_{\B} (b_m))_m\subseteq T\B$,\footnote{The use of order limits here is motivated by the fact that even when $(a_n)$ and $(b_m)$ are monotone increasing, $a_n \otimes b_m$ may not be. For example, if $b_m$ is negative, then $(a_n \otimes b_m)_n$ is monotone decreasing.} 
% \begin{equation}\label{eq:theta_description}
%   \theta_{\A,\B}(
%   \olim_n^{T\A}\eta_{\A}(a_n)\otimes
%   \olim_m^{T\B}\eta_{\B}(b_m)) =  \olim_n^{T(\A\otimes\B)}
%       \olim_m^{T(\A\otimes\B)} \bigl(\eta_{\A\otimes\B}(a_n\otimes b_m)\bigr)
% \end{equation}
% (cf. \cref{def:bisigmanormal}).

\begin{remark}\label{rem:monoidalmonad_cpu}
	The extension $\bar T = \bar U \bar F\colon \Calgcpu \to \Calgcpu$ is also a symmetric monoidal monad by the same reasoning of \cref{prop:monoidalmonad}, since the Kleisli category of $\bar T$ is given by Pedersen--Baire envelopes with $\sigma$-cpu maps, and this was shown to be a symmetric monoidal category in \cite[Remark 5.3.3]{fritz2025categories}.
\end{remark}

To be sufficiently self-contained, we now provide some details on the categorical results later applied to ensure that $\SigmaC$ and $\SigmaCcpu$ are symmetric monoidal categories.
Although the main reference~\cite{seal2013tensors} deals with monoidal monads and monoidal categories, here we restrict to \emph{symmetric} monoidal structures, as the additional generality never surfaces in our investigation.
\begin{notation} We adopt the following notation.
	\begin{itemize}
		\item The category of Eilenberg--Moore algebras of a monad $T\colon \clK\to \clK$ is denoted by $\EM(T)$ and its objects are denoted by pairs $(A,a)$ where $\A\in \clK$ and $a\colon T\A \to \A$ denotes the algebra structure.
		\item A symmetric monoidal category is a triple $(\clK, \otimes , I)$ where $\clK$ is a category, $\otimes \colon \clK \times \clK \to \clK$ denotes the tensor (monoidal) product, and $I$ denotes the unit: $\A \otimes I \cong \A$ for all $\A \in \clK$.
		\item A symmetric monoidal monad is a pair $(T,\theta)$ where $T$ is a monad on a symmetric monoidal category $(\clK,\otimes ,I)$, and $\theta$ is a natural transformation given by $\theta_{\A,\B} \colon T\A \otimes T\B \to T(\A\otimes \B)$.
	\end{itemize}
\end{notation}


\begin{definition}\label{def:kocktensor}
	Let $(T,\theta)$ be a symmetric monoidal monad on a symmetric monoidal category $(\clK,\otimes,I)$.
	The \newterm{tensor product} of Eilenberg--Moore algebras $(\A,a)$ and $(\B,b)$ is given by the coequalizer diagram
	\begin{equation*}%\label{eq:kocktensor}
		\xymatrix@C=5em{
		T(T\A\otimes T\B)
		\ar@<.7ex>[r]^{\mu T\theta }
		\ar@<-.7ex>[r]_{T(a\otimes b)}
		& T(\A\otimes \B)\ar[r]^{q_{(\A,a),(\B,b)}} & (\A,a)\kocktensor (\B,b)
		}
	\end{equation*}
	in $\EM(T)$.
	When the algebra morphisms are clear from context (or irrelevant), we simply write $\A \kocktensor \B$ and $q_{\A,\B}$.

	We say that $\EM(T)$ \newterm{has all tensors} if $q_{\A,\B}$ and $\A \kocktensor \B$ exist for all $(\A,a)$ and $(\B,b)$ in $\EM(T)$.
\end{definition}
\begin{remark}\label{rem:coeq_homtensor}
	\begin{itemize}
		\item  The coequalizer in the definition is reflexive with common section $T(\eta_A\otimes \eta_B)$.
		\item Although additional requests are needed to ensure one has a symmetric monoidal category (e.g., $\otimes$ should be associative), the above definition is sufficient to ensure the existence of $\phi\kocktensor \psi$ for morphisms $\phi\colon (\A,a) \to (\C,c)$ and $\psi\colon (\B,b)\to (\D,d)$, because of coequalizer properties:
		      \[
			      \xymatrix@C=5em{
			      T(\A\otimes \B)\ar[r]^{q_{\A,\B}}\ar[d]_{T(\phi\otimes \psi)} & \A\kocktensor \B\ar@{-->}[d]^{\phi\kocktensor \psi} \\
			      T(\C\otimes \D)\ar[r]^{q_{\C,\D}} & \C\kocktensor \D
			      }
		      \]
	\end{itemize}

\end{remark}



\begin{proposition}[{\cite[Lemma 2.3.3]{seal2013tensors}}]\label{prop:bijection_bimorphisms} Let $(T,\theta)$ be a symmetric monoidal monad on a symmetric monoidal category $(\clK,\otimes,I)$ such that $\EM(T)$ has all tensors.
	Consider $(\A,a)$, $(\B,b)$ and $(\C,c)$ Eilenberg--Moore algebras.
	The morphism
	\[
		\xymatrixcolsep{3pc}\xymatrix{
		\A \otimes \B \ar[r]^(0.45){\eta_{\A\otimes \B}}
		& T(\A\otimes \B) \ar[r]^-{q_{\A,\B}}
		& \A\kocktensor\B
		}
	\]
	induces, by precomposition, a bijection between the following:
	\begin{itemize}
		\item Morphisms of type $\A \kocktensor \B \to \C$.
		\item Bimorphisms of $\EM(T)$, i.e.\ morphisms $\omega\colon \A\otimes \B\to \C$ in $\clK$ satisfying
		      \[
			      \xymatrixcolsep{3pc}\xymatrix{
			      T\A \otimes T\B \ar[r]^{\theta_{\A,\B}} \ar[d]_{a\otimes b}
			      & T(\A\otimes \B) \ar[r]^-{T\omega}
			      & T\C \ar[d]^{c} \\
			      \A\otimes \B \ar[rr]_{\omega}
			      & & \C
			      }
		      \]
	\end{itemize}
\end{proposition}
\begin{theorem}[{\cite[Theorem 2.5.5]{seal2013tensors}}]\label{thm:SMC1}
	Let $(T,\theta)$ be a symmetric monoidal monad on a symmetric monoidal category $(\clK,\otimes,I)$ such that $\EM(T)$ has all tensors.
	If for all Eilenberg--Moore algebras $(\A,a)$, $(\B,b)$, $(\C,c)$, and $(\D,d)$,
	\begin{enumerate}
		\item \label{it:coeq_associators}
		      The diagrams of the form
		      \begin{equation*}
			      \begin{array}{c}
				      \xymatrix@C=6em{
				      T(T\A\otimes T\B)\kocktensor TT\C
				      \ar@<.7ex>[r]^{\mu T\theta \kocktensor \mu}
				      \ar@<-.7ex>[r]_{T(a\otimes b)\kocktensor Tc}
				       & T(\A\otimes \B)\kocktensor T\C \ar[r]^{q\kocktensor c}  & (\A\kocktensor \B)\kocktensor \C
				      }                                                                                             \\
				      \xymatrix@C=6em{
				      TT\A \kocktensor T(T\B\otimes T\C)
				      \ar@<.7ex>[r]^{\mu  \kocktensor \mu T\theta }
				      \ar@<-.7ex>[r]_{Ta \kocktensor T(b\otimes c)}
				       & T\A \kocktensor T(\B\otimes \C) \ar[r]^{a\kocktensor q} & \A\kocktensor (\B\kocktensor \C)
				      }
			      \end{array}
		      \end{equation*}
		      are coequalizers in $\EM(T)$;
		\item \label{it:coeq_tensor_epic}
		      $q_{\A,\B}\kocktensor q_{\C,\D}$ are epimorphisms in $\EM(T)$;
	\end{enumerate}
	Then $(\EM(T),\kocktensor,TI)$ is a symmetric monoidal category.
\end{theorem}

\begin{theorem}[{\cite[Theorem 2.6.4]{seal2013tensors}}]\label{thm:SMC2}
	Let $(T,\theta)$ be a symmetric monoidal monad on a symmetric monoidal category $(\clK,\otimes ,I)$ such that $\EM(T)$ has all tensors.
	If $T(\A\otimes -)$ preserves reflexive coequalizers in $\clK$ for all $\A \in \clK$, then $(\EM(T),\kocktensor,TI)$ is a symmetric monoidal category.
\end{theorem}
As it will be important in the proof of \cref{thm:tensor_scpu}, we mention that the proof of \cref{thm:SMC2} consists in showing that the assumptions of \cref{thm:SMC1} are satisfied.

\begin{proposition}[{\cite[Theorem 3.3]{khoshkam2002categorical}}]\label{prop:tensor_colims}
	Let $\A\in \Calg$.
	Maximal tensoring $\A\otimes - \colon \Calg\to \Calg$ preserves connected colimits.
\end{proposition}
%This fact seems somehow known, although it is hard to find an explicit proof.\footnote{One considering inductive limits can be found at \url{https://mathoverflow.net/questions/381936}.} 
For completeness, we provide a direct reasoning using the universal property of the maximal tensor product.
\begin{proof}
	Let us consider a connected category $\clC$ and a functor $L:\clC\to \Calg$ admitting a colimit $\colim L$.
	We now consider a cocone $\phi_c \colon Lc\otimes \A \to \B$ of $L(-)\otimes \A$ for some $\B\in \Calg$.
	We seek to show that there exists a compatible map $\colim L \otimes \A\to \B$, thereby ensuring that $\colim(L(-)\otimes A)\cong (\colim L)\otimes A$.

	Since $\clC$ is connected, one has a zig-zag of morphisms from any $c$ to $c'$.
	For simplicity, let us consider the case $f\colon c\to c'$.
	The commutative diagram
	\[
		\xymatrix{
		Lc\otimes \A \ar[rr]^{Lf \otimes \id_{\A}} \ar[dr]_{\phi_c}
		&&
		Lc'\otimes \A \ar[dl]^{\phi_{c'}}
		\\
		&
		\B
		&
		}
	\]
	ensures that $\phi_c(1\otimes a) = \phi_{c'}(1\otimes a)$.
	Since $\clC$ is connected, we obtain a unique map $\psi\colon \A \to \B$, which sends $a$ to $\phi_{c}(1\otimes a)$ for any choice of $c$.
	By looking at the marginalization on the other side, $\phi_c(-\otimes 1)\colon Lc\to \B$ induce a cocone for $L$. Therefore, one obtains $\phi\colon \colim L \to \B$.

	We now aim to show that the image of $\phi$ and $\psi$ commute. Once this is established, the universal property of the maximal tensor product would immediately yield $\colim L \otimes \A \to \B$, thus concluding the proof.

	Let us consider $x\in \colim L$. If $x$ is in the image of some $Lc$, then $\phi(x)$ commutes with the image of $\psi$ by construction.
	If not, $x$ is a uniform limit of a sequence of elements commuting with the image of $\psi$, because colimits in $\Calg$ are given by the completion of the colimits of $*$-algebras over $\bbC$ (see \cite[Corollaries 7.22 and 7.23]{bunke2021lectures}).\footnote{This fact was already used in \cref{lem:technical_access}.}
	In this case, one notes that $\phi(x)\psi(a)$ and $\psi(a)\phi(x)$ are both limits of $\phi(x_n)\psi(a)=\psi(a)\phi(x_n)$, and therefore they coincide.
	Finally, $\phi$ and $\psi$ commute, and the statement follows.
\end{proof}

\begin{theorem}\label{thm:tensor_shom}
	$\SigmaC$ is a symmetric monoidal category whose monoidal unit is $\bbC$ and whose tensor product $\kocktensor$ satisfies:
	\begin{enumerate}
		\item\label{it:free_algebras} $F(\A)\kocktensor F(\B)\cong F(\A\otimes \B)$;
		\item\label{it:universal_prop} The following universal property:
		      for every $\phi \colon \A \to \C$ and $\psi\colon \B \to \C$ in $\SigmaC$ with commuting images, there exists a unique $(\phi,\psi)\colon \A \kocktensor \B \to \C$ satisfying $(\phi,\psi)(a\kocktensor b) = \phi(a)\psi(b)$.
		      Moreover, all \starshom{}s of type $\A \kocktensor \B \to \C$ are obtained in this way.
	\end{enumerate}
\end{theorem}
\begin{proof}
	From \cref{prop:cocomplete,prop:tensor_colims}, $T(\A \otimes -)$ preserve reflexive coequalizers, and moreover $\SigmaC$ has tensors because it is cocomplete.
	By \cref{thm:SMC2}, $\SigmaC$ is a symmetric monoidal category (recall that $T$ is a symmetric monoidal monad by \cref{prop:monoidalmonad}).
	The monoidal unit is $\bbC$ because it is isomorphic to $T(\bbC)$ (for every finite-dimensional \cstar{} $\A$, $T(\A)\cong \A$, see \cite[Example 3.3.16]{fritz2025categories}).
	Item~\ref{it:free_algebras} follows from \cite[2.5.3]{seal2013tensors}.

	We are therefore reduced to show Item~\ref{it:universal_prop}.
	By \cref{prop:bijection_bimorphisms}, the morphisms $\A \kocktensor \B \to \C$ are in bijective correspondence with the bimorphisms $\A\otimes \B\to\C$ of $\EM(T)$ via precomposing with $q_{\A,\B}\eta_{\A\otimes \B}$.
	We claim that such morphisms exactly correspond to a choice of $\phi$ and $\psi$ as stated in Item~\ref{it:universal_prop}, which will then conclude the proof.

	Starting with a bimorphism $\omega\colon \A\otimes \B\to \C$, $\omega$ is uniquely determined by its marginals $\phi=\omega(-\otimes 1)$ and $\psi=\omega(1\otimes -)$ by the universal property of the maximal tensor product.
	Moreover, $\phi$ and $\psi$ have commuting images by direct check:
	\begin{equation*}
		\begin{split}
			\phi(a)\psi(b)&= \omega(a\otimes 1)\omega(1\otimes b) = \omega(a\otimes b) \\
			&= \omega((1\otimes b)(a\otimes 1))= \omega(1\otimes b)\omega(a\otimes 1) \\
			&= \psi(b)\phi(a).
		\end{split}
	\end{equation*}
	It then suffices to ensure that $\phi$ and $\psi$ so defined are necessarily $\sigma$-normal.
	For $\phi$, this follows by the commutative diagram
	\[
		\xymatrix@C=4em{
			\A\ar[r]^(0.45){\cong}\ar[rd]_{\cong}\ar@/^2em/[rrr]^{\phi} & \A \otimes \bbC \ar[r]^{\id_{\A}\otimes \, !}\ar[d]^{q\eta} & \A\otimes \B \ar[r]^-{\omega}\ar[d]^{q\eta}& \C\\
		&\A \kocktensor \bbC \ar[r]^{\id_{\A}\kocktensor \, !} & \A\kocktensor \B \ar[ru]&
		}
	\]
	since the maps in the lower path are \starshom{}s (where $!$ denotes the unique unital linear map from $\bbC$). Similarly, also $\psi$ is a \starshom.

	Conversely, let us consider two $\sigma$-homomorphisms $\phi$ and $\psi$ as in the statement.
	These yield a map $\omega \colon \A \otimes \B \to \C$ given by $\omega (a\otimes b)=\phi(a)\psi(b)$ via the universal property of the maximal tensor product. We claim that such an $\omega$ is a bimorphism in the sense of \cref{prop:bijection_bimorphisms}.
	By again employing the universal property of the maximal tensor product, the two compositions $T\A \otimes T\B\to \C$ can be tested on their decompositions $T\A \to \C$ and $T\B \to \C$.
	Since these are both $\sigma$-homomorphisms and $\C$ is a \cstars{}, both are uniquely determined by the $*$-homomorphisms $\A \to \C$ and $\B \to \C$ given by composing with the unit $\eta$.
	It is now a simple check that $\phi$ corresponds to both
	\[
		\A \to T(\A\otimes \B) \to T\C \to \C \qquad \text{and}\qquad \A \to \A\otimes \B\to \C
	\]
	and analogously for $\psi$, so the wanted commutation holds.
\end{proof}


\begin{proposition}\label{prop:tensor_commutative}
	Given $\A$, $\B$ commutative \cstars{}s, their tensor product $\A \kocktensor \B$ is again commutative.
	Moreover, $\A\kocktensor \B$ is the coproduct in $\cSigmaC$.
\end{proposition}

\begin{proof}
	Whenever $\A$ and $\B$ are commutative, $\A \otimes\B$ is commutative  (see, e.g., \cite[Lemma T.6.16]{weggeolsen1993Ktheory}).
	Moreover, $T$ preserves commutativity~\cite[Remark 5.3.4(iii)]{saito15monotone}, hence $T(\A\otimes \B)$ is also commutative.
	Since commutativity is preserved by quotients and $\A \kocktensor \B$ is a quotient of $T(\A\otimes \B)$ (by definition of the tensor and \cref{lem:reflexive_coeq}), we obtain that $\A\kocktensor \B$ is commutative.
	The second part of the statement is a direct application of the universal property of the tensor product (\cref{thm:tensor_shom}\ref{it:universal_prop}).
\end{proof}

\begin{remark}\label{rem:tensor_bool}
	Recall that $\cSigmaC$ is equivalent to the category of \bsigma{}s and $\sigma$-homomorphisms, as mentioned in \cref{sec:preliminaries}.
	By combining this equivalence with \cref{prop:tensor_commutative}, $\kocktensor$ retrieves the universal tensor product of \bsigma{}s \cite[Proposition 2.5.1]{fritz2025categories}.

	In particular, given two measurable spaces $X$ and $Y$, it is not always true that $\Linf(X)\kocktensor \Linf(Y)\cong \Linf(X\times Y)$ (cf.~\cite[Example 2.5.6]{fritz2025categories}, or the historical reference \cite[Example 37.A]{sikorski69boolean}).
	However, if $X$ and $Y$ are Baire measurable spaces (this notion was recalled before \cref{thm:giry_restriction}), we have that $\Linf(X)\cong F(C(X))$, and therefore
	\begin{align*}
		\Linf(X)\kocktensor \Linf(Y) &\cong F(C(X)) \kocktensor F(C(Y)) \\
		& \cong F(C(X)\otimes C(Y)) \\
		& \cong F(C(X\times Y)) \\
		& \cong \Linf(X\times Y)
	\end{align*}
	where $C(X)\otimes C(Y)\cong C(X\times Y)$ is well-known (see, e.g., \cite[Lemma T.6.16]{weggeolsen1993Ktheory}).
\end{remark}
% \begin{remark}
%   An alternative proof of \cref{prop:tensor_commutative}, suggested to us by Tobias Fritz, consists in noting that $\A^{\op}\kocktensor \B^{\op}\cong (\A\kocktensor \B)^{\op}$, where $-^{\op}$ denotes the opposite \cstar{} (i.e.\ where multiplication is reversed, $a\cdot^{\op} b\coloneqq a\cdot b$).
%   Indeed, the isomorphisms $\A^{\op} \otimes \B^{\op} \cong( \A\otimes \B)^{\op}$ and $T(\A^{\op})\cong T(\A)^{\op}$ are easily checked with sufficient grasp of the concepts. 
% \end{remark}

We now extend the tensor product to $\SigmaCcpu$.
As we will see, the proof boils down to noting that the coequalizers needed are reflexive and live in $\SigmaC$, and so \cref{lem:reflexive_coeq} applies.
\begin{theorem}\label{thm:tensor_scpu}
	$\SigmaCcpu$ is a symmetric monoidal category by extending the tensor product $\kocktensor$ of $\SigmaC$.
\end{theorem}
\begin{proof}
	We check that the hypotheses of \cref{thm:SMC1} are satisfied. Recall that $\bar{T}$ is a symmetric monoidal monad (\cref{rem:monoidalmonad_cpu}).

	We note that the algebra structures are uniquely determined, namely $a=U\epsilon_{\A}$, $b=U\epsilon_{\B}$ and $c=U\epsilon_{\C}$, as discussed at the end of the proof of \cref{thm:U_monadic}.
	In particular, the algebra structures are \starshom{}s, i.e.~they are in the image of the forgetful functor $\SigmaC\to \SigmaCcpu$.
	Since this functor preserves reflexive coequalizers by \cref{lem:reflexive_coeq}, we obtain that $\SigmaCcpu$ has tensors (the fact that tensors are reflexive coequalizer was noted in \cref{rem:coeq_homtensor}).

	To ensure that the diagrams of Item~\ref{it:coeq_associators} in \cref{thm:SMC1} are coequalizers, we first note that they already are so in $\SigmaC$ because \cref{thm:tensor_shom} applies \cref{thm:SMC2}, and the proof of this result consists in checking that the hypotheses of \cref{thm:SMC1} hold.
	Moreover, these coequalizer diagrams are reflexive: the common sections are $T(\eta_{\A}\otimes \eta_{\B})\kocktensor T(\eta_{\C})$ and $T(\eta_{\A})\kocktensor T(\eta_{\B}\otimes \eta_{\C})$, respectively.
	For the first coequalizer in Item~\ref{it:coeq_associators}, using functoriality of $T$ and $\kocktensor$, we have
		\begin{align*}
			(\mu T\theta \kocktensor \mu )(T(\eta_{\A}\otimes \eta_{\B})\kocktensor T(\eta_{\C}) ) & = \mu T(\theta (\eta_{\A}\otimes \eta_{\B}))\kocktensor \mu T(\eta_{\C})\\
			                                                                                       & = \mu T(\eta_{\A\otimes \B}) \kocktensor \id_{T\C}                      \\
			                                                                                       & = \id_{T(A\otimes B)}\kocktensor \id_{T\C}                              \\
			                                                                                       & = \id_{T(\A\otimes B)\kocktensor T\C}                                   \\
			(T(a\otimes b)\kocktensor Tc)(T(\eta_{\A}\otimes \eta_{\B})\kocktensor T(\eta_{\C}) )  & = T(a\eta_{\A}\otimes b\eta_{\B}) \kocktensor T(c\eta_{\C})             \\
			                                                                                       & = T(\id_{\A\otimes \B})\kocktensor T(\id_{\C})                          \\
			                                                                                       & = \id_{T(\A\otimes B)\kocktensor T\C}.
		\end{align*}
	An analogous reasoning applies to the second coequalizer in Item~\ref{it:coeq_associators}.
	As $\SigmaC \to \SigmaCcpu$ preserves reflexive coequalizers (\cref{lem:reflexive_coeq}), Item~\ref{it:coeq_associators} is verified.

	We now discuss Item~\ref{it:coeq_tensor_epic} of \cref{thm:SMC1}.
	We claim that $q_{\A,\B}\kocktensor q_{\C,\D}$ are in fact surjective maps.
	We first note that $q_{\A,\B}$ are all surjective by \cref{lem:reflexive_coeq}, and moreover there is a commutative diagram
	\[
		\xymatrix@C=5em{
		T(T(\A\otimes \B) \otimes T(\C\otimes \D))
		\ar[r]^{T(q\otimes q)}
		\ar@{->>}[d]_{q}
		&
		T(\A\kocktensor \B \otimes \C\kocktensor \D)
		\ar@{->>}[d]^{q}
		\\
		T(\A\otimes \B) \kocktensor T(\C\otimes \D)
		\ar[r]_{q\kocktensor q}
		&
		(\A\kocktensor \B) \kocktensor (\C\kocktensor \D)
		}
	\]
	In particular, the vertical arrows are surjective as already noted.
	To conclude, we only need to show that $T(q\otimes q)$ is also surjective. This follows from the fact that the maximal tensor product is exact~\cite[Proposition 3.7.1]{brownozawa} and $T$ preserves surjectivity (\cref{rem:T_surjective}).
\end{proof}


We now aim to describe all the \scpu{} maps of type $\A\kocktensor \B\to \C$.
To this end, we give the following.
\begin{definition}\label{def:bisigmanormal}
	Let us consider three \cstars{} $\A$, $\B$ and $\C$. A \cpu{} map $\omega\colon \A \otimes \B \to \C$ is \newterm{bi-$\sigma$-normal} if
	\[
		\omega\left( \olim_n^{\A} a_n \otimes \olim_m^{\B} b_m \right) = \olim_n^{\C} \olim_m^{\C} \omega(a_n \otimes b_m)
	\]
	for all order convergent sequences $(a_n)\subseteq \A$ and $(b_m)\subseteq \B$.
\end{definition}

\begin{lemma}\label{lem:theta_bisigma}
	The map $\theta\colon T\A \otimes T\B \to T(\A\otimes \B)$ is bi-$\sigma$-normal.
\end{lemma}
\begin{proof}
	By construction, $\theta$ is a $*$-homomorphism, so $\theta(a\otimes b)=\theta(a\otimes 1)\theta(1\otimes b)$, and $\theta(-\otimes 1)\colon T\A \to T(\A\otimes \B)$ is a $\sigma$-homomorphism, so $\theta(\olim_n a_n \otimes 1) = \olim_n \theta(a_n \otimes 1)$, and similarly for the second argument.

	Consequently, the statement holds by compatibility of order limits with multiplication:
	\begin{equation*}
		\begin{split}
			\theta\left(\olim_n^{T\A} a_n \otimes \olim^{T\B}_m b_m\right) & =\theta\left(\olim^{T\A}_n a_n \otimes 1\right)\theta\left(1\otimes \olim^{T\B}_m b_m\right)    \\
			                                                               & = \olim^{T(\A\otimes \B)}_n \theta(a_n \otimes 1)\olim^{T(\A\otimes \B)}_m \theta(1\otimes b_m) \\
			                                                               & = \olim^{T(\A\otimes \B)}_n \olim^{T(\A\otimes \B)}_m \theta(a_n \otimes 1)\theta(1\otimes b_m) \\
			                                                               & =\olim^{T(\A\otimes \B)}_n \olim^{T(\A\otimes \B)}_m \theta(a_n \otimes b_m).
		\end{split}
	\end{equation*}
\end{proof}
\begin{proposition}[Universal property of $\kocktensor$]\label{prop:universal_bisigma}
	The morphisms $\A \kocktensor \B \to \C$ in $\SigmaCcpu$ are in bijective correspondence to those $\omega\colon \A \otimes \B \to \C$ in $\Calgcpu$ that are bi-$\sigma$-normal.
\end{proposition}
\begin{proof}
	It suffices to show that $\omega\colon \A \otimes \B\to \C$ is bi-$\sigma$-normal if and only if is a bimorphism in the sense of \cref{prop:bijection_bimorphisms}.
	As already mentioned in the proof of \cref{thm:tensor_scpu}, the algebra structures are uniquely determined as shown at the end of the proof of \cref{thm:U_monadic}.
	We can therefore assume $a=U\epsilon_{\A}$, $b=U\epsilon_{\B}$ and $c=U\epsilon_{\C}$ in the bimorphism diagram of \cref{prop:bijection_bimorphisms}.

	If $\omega$ is a bimorphism, we take order convergent sequences $(a_n)\subseteq \A $ and $(b_m)\subseteq \B$.
	We wish to unravel the commuting diagram describing bimorphisms on the element $\olim_n \eta_{\A}(a_n)\otimes \olim_m \eta_{\B}(b_m)$.

	Concerning the upper path, we have
	\begin{align*}
		U\epsilon_{\C} T\omega \theta_{\A,\B} (\olim_n \eta_{\A}(a_n)\otimes \olim_m \eta_{\B}(b_m)) & = U\epsilon_{\C} T\omega (\olim_n\olim_m \theta_{\A,\B} (\eta_{\A}(a_n)\otimes \eta_{\B}(b_m)))  \\
		&= \olim_{n}\olim_m  U\epsilon_{\C} T\omega \theta_{\A,\B} (\eta_{\A}(a_n)\otimes \eta_{\B}(b_m))
	\end{align*}
	where we used \cref{lem:theta_bisigma} and $\sigma$-normality of $T\omega$ and $U\epsilon_{\C}$.
	Moreover, the following diagram commutes by unit compatibility of $\theta$ and naturality of $\eta$:
	\begin{equation}\label{eq:omega_comm}
		\xymatrixcolsep{3pc}\xymatrix{
		&\A\otimes \B\ar[ld]_{\eta_{\A}\otimes \eta_{\B}} \ar[d]^{\eta_{\A\otimes \B}} \ar[r]^{\omega}& \C\ar[d]^{\eta_{\C}}\ar[rd]^{\id_{\C}}\\
		T\A \otimes T\B \ar[r]^{\theta_{\A,\B}}
		& T(\A\otimes \B) \ar[r]^-{T\omega}
		& T\C \ar[r]^{U\epsilon_{\C}} & \C
		}
	\end{equation}
	Putting the two together, the upper path yields $\olim_n \olim_m \omega (a_n \otimes b_m)$.

	For the lower path,
	\[
		\omega (U(\epsilon_{\A})\otimes U(\epsilon_{\B} )(\olim_n \eta_{\A}(a_n)\otimes \olim_m \eta_{\B}(b_m)) = \omega (\olim_n a_n \otimes \olim_m b_m)
	\]
	by $\sigma$-normality of $U(\epsilon_{\A})$ and $U(\epsilon_{\B})$.
	Finally, by commutativity of the bimorphism diagram in \cref{prop:bijection_bimorphisms},
	\[
		\omega (\olim_n a_n \otimes \olim_m b_m) = \olim_n \olim_m \omega (a_n \otimes b_m),
	\]
	i.e.\ $\omega$ is bi-$\sigma$-normal.

	We now seek to prove the converse, i.e.\ that any bi-$\sigma$-normal ${\omega}$ is a bimorphism.
	This is done by two $\pi$-$\lambda$ arguments.
	For the first one, let
	\[
		M_{\B} \coloneqq \lbrace a \in T\A \mid \omega(U\epsilon_{\A}\otimes U\epsilon_{\B})(a \otimes \eta_{\B}(b)) = U\epsilon_{\C} T\omega \theta_{\A,\B}(a \otimes \eta_{\B}(b)) \text{ for all }b\in\B \rbrace\subseteq T\A
	\]
	This contains $\eta_{\A}(\A)$, because precomposing with $\eta_{\A}\otimes \id_{\B}$ yields $\omega$ on both sides of the defining equation by commutativity of \eqref{eq:omega_comm}.
	Moreover, $M_{\B}$ is a $*$-subspace because all maps involved are $*$-preserving and linear.
	Next, if $(a_n)\subseteq T\A$ is a bounded monotone increasing sequence contained in $M_{\B}$, then $\sup_n^{T\A} (a_n)$ belongs to $M_{\B}$ by bi-$\sigma$-normality of $\omega$ and $\theta$ (established in \cref{lem:theta_bisigma}).
	Since $T\A$ is the $\sigma$-closure of $\A$, we conclude that $M_{\B}=T\A$.

	For the second $\pi$-$\lambda$ argument, consider the $*$-subspace
	\[
		M_{T\A} \coloneqq \lbrace b \in T\B \mid \omega(U\epsilon_{\A}\otimes U\epsilon_{\B})(a \otimes b) = U\epsilon_{\C} T\omega\theta_{\A,\B}(a \otimes b) \text{ for all }a\in T\A\rbrace\subseteq T\B
	\]
	By the previous discussion, $\eta_{\B}(\B)\subseteq M_{T\A}$.
	Moreover, the argument above can be adapted to show that $M_{T\A}$ also is a $\sigma$-subspace, so $M_{T\A}=T\B$.
	Finally, the set
	\[
		M\coloneqq \lbrace z \in T\A\otimes T\B \mid \omega(U\epsilon_{\A}\otimes U\epsilon_{\B})(z) = U\epsilon_{\C} T\omega\theta_{\A,\B}(z)\rbrace\subseteq T\A\otimes T\B
	\]
	contains all simple tensors $z=a\otimes b$.
	$M$ is a norm-closed $*$-subspace, where norm-closure follows from the fact that all the maps are contractive.
	Therefore, $M$ contains all the norm limits of linear combinations of simple tensors, hence $M=T\A\otimes T\B$.
\end{proof}


% \begin{remark}
%   %This remark requires a hierarchical reasoning... Maybe not that important here, also the machinery would have to be introduced exactly for this remark, doesn't seem worth it.
%   A more explicit description of $\A \kocktensor \B$ is given by taking $T(\A\otimes \B)$ and quotienting it via the $\sigma$-ideal generated by
%   \[
%   \Im (\delta - \eta_{\A \otimes \B}(U\epsilon_{\A}\otimes U\epsilon_{\B}))
%   \]
%   This is an immediate consequence of the fact that the $\sigma$-ideal considered to define the tensor product is $\Im(\mu T\delta - T(U\epsilon_{\A}\otimes U\epsilon_{\B}))$, see \eqref{eq:tensor_def} in the proof of \cref{thm:tensor_scpu}.
%   Since $\mu T\delta $ and $T(U\epsilon_{\A}\otimes U\epsilon_{\B})$ are $\sigma$-homomorphisms, they are uniquely determined by their composition with the unit $\eta_{T\A\otimes T\B}$, from which one obtains
%   \[
%   \mu T(\delta) \eta_{T\A\otimes T\B} = \mu \eta_{T(\A\otimes \B)} \delta = \delta \qquad\text{ and }\qquad T(U\epsilon_{\A}\otimes U\epsilon_{\B})\eta_{T\A\otimes T\B} = \eta_{\A\otimes \B} (U\epsilon_{\A}\otimes U\epsilon_{\B}),
%   \] 
%   where both equalities follow from naturality of $\eta$.
% \end{remark}

